\documentclass[12pt]{article}
\usepackage[utf8]{inputenc}
\usepackage[spanish]{babel}
\usepackage[a4paper, margin=2cm]{geometry}
\usepackage{tikz}
\usetikzlibrary{babel}
\usepackage[most]{tcolorbox}
\usepackage{amsmath}

% Usar fuente sin serifa para un aspecto más moderno (tipo infografía)
\renewcommand{\familydefault}{\sfdefault}

% Definición de colores con temática médica/científica
\definecolor{medblue}{RGB}{43, 108, 176}
\definecolor{medteal}{RGB}{49, 151, 149}
\definecolor{medlight}{RGB}{235, 248, 255}
\definecolor{meddark}{RGB}{26, 32, 44}
\definecolor{bloodred}{RGB}{220, 38, 38}   % Rojo para la sangre
\definecolor{aortared}{RGB}{239, 68, 68}    % Rojo brillante para la arteria

\begin{document}
\pagestyle{empty}

\begin{tcolorbox}[
    enhanced, 
    colback=medlight!30, 
    colframe=medblue, 
    title={\Large \textbf{Sistemas de Unidades: Velocidad de Flujo}}, 
    halign title=center, 
    fonttitle=\bfseries, 
    arc=4mm, 
    boxrule=1.5pt,
    shadow={2mm}{-2mm}{0mm}{black!30!white}
]

    \vspace{0.2cm}
    \begin{tcolorbox}[colback=white, colframe=medteal, arc=2mm, boxrule=1.2pt, title=\textbf{Caso Clínico / Problema}]
        La velocidad media del flujo sanguíneo en la aorta humana es de aproximadamente \textbf{30 cm/s}. Convierta esta velocidad a metros por hora (\textbf{m/h}).
    \end{tcolorbox}

    \vspace{0.4cm}
    
    \begin{center}
    \begin{tikzpicture}[scale=0.95]
        
        % Vaso sanguíneo 3D (Aorta)
        % Cara posterior y cuerpo
        \filldraw[aortared!30, draw=bloodred, thick] (-4,-1.2) -- (4,-1.2) arc (-90:90:0.5cm and 1.2cm) -- (-4,1.2) arc (90:-90:0.5cm and 1.2cm);
        % Cara frontal (apertura izquierda)
        \filldraw[aortared!60, draw=bloodred, thick] (-4,0) ellipse (0.5cm and 1.2cm);
        
        % Líneas de flujo internas
        \draw[->, line width=3pt, white] (-3.5, 0) -- (3.5, 0) node[midway, above, font=\bfseries, text=white] {Flujo Sanguíneo};
        \draw[->, thick, white!80!aortared] (-3.2, 0.6) -- (2.5, 0.6);
        \draw[->, thick, white!80!aortared] (-3.2, -0.6) -- (2.5, -0.6);

        % Etiquetas de velocidad arriba
        \node[font=\Large\bfseries, text=bloodred, fill=white, inner sep=4pt, rounded corners=3pt, draw=bloodred!30, thick] at (-3, 2.3) {$30 \text{ cm/s}$};
        \node[font=\Large\bfseries, text=meddark, fill=white, inner sep=4pt, rounded corners=3pt, draw=meddark!30, thick] at (3, 2.3) {$1080 \text{ m/h}$};

        % Flecha de conversión arriba
        \draw[->, ultra thick, medteal] (-1.3, 2.3) -- (1.3, 2.3) node[midway, above, font=\small\bfseries] {Conversión};

        % Factores de conversión debajo
        \node[font=\Large, text=meddark] at (0, -2.8) {
            $v = 30 \, \frac{\text{cm}}{\text{s}} \times \underbrace{\left( \frac{1 \text{ m}}{100 \text{ cm}} \right)}_{\text{Distancia}} \times \underbrace{\left( \frac{3600 \text{ s}}{1 \text{ h}} \right)}_{\text{Tiempo}}$
        };
        
    \end{tikzpicture}
    \end{center}

    \vspace{0.4cm}
    \begin{tcolorbox}[colback=medteal!10, colframe=medteal, arc=2mm, boxrule=1.2pt, title=\textbf{Desarrollo Matemático y Solución}]
        Para convertir unidades compuestas de velocidad, utilizamos factores de conversión tanto para la longitud ($1 \text{ m} = 100 \text{ cm}$) como para el tiempo ($1 \text{ h} = 3600 \text{ s}$):
        \begin{align*}
            v &= 30 \, \frac{\text{cm}}{\text{s}} \times \left( \frac{1 \text{ m}}{100 \text{ cm}} \right) \times \left( \frac{3600 \text{ s}}{1 \text{ h}} \right) \\[1.5ex]
            v &= 30 \times 0.01 \times 3600 \, \frac{\text{m}}{\text{h}} = \mathbf{1080 \, \frac{\text{m}}{\text{h}}}
        \end{align*}
        \textbf{Resultado:} La velocidad media del flujo sanguíneo en la aorta equivale a \textbf{1080 m/h}.
    \end{tcolorbox}
    
\end{tcolorbox}

\end{document}